\documentclass[10pt]{article}
\usepackage[T1]{fontenc}
\usepackage[utf8]{inputenc}
\usepackage[french]{babel}
\usepackage{amsmath,amssymb,mathrsfs,mathpazo}
\usepackage{geometry}
\geometry{a4paper,margin=2cm}
\usepackage{graphicx,array,tabularx,tabularray,colortbl,multirow,multicol,booktabs,enumitem,pifont,ulem,fancybox,diagbox,cancel}
\usepackage{pstricks,pstricks-add,pst-plot,pst-tree,pst-node,pst-text,pst-func,pst-math,pst-bspline,pst-3dplot,pst-eucl,pst-solides3d}
\let\pscalculate\undefined
\usepackage{pst-bezier}
\usepackage{tikz,pgfplots,tkz-tab}
\usetikzlibrary{arrows,arrows.meta,patterns,calc,positioning,trees,decorations.markings,shapes.geometric,shapes.misc,shapes.arrows,intersections,decorations.pathmorphing,decorations.pathreplacing,backgrounds,shadows,babel}
\pgfplotsset{compat=1.18}
\usepgfplotslibrary{fillbetween}
\usepackage[np]{numprint}
\usepackage{listings,makecell,pas-tableur,twoopt}
\let\barre\undefined
\usepackage[pstricks]{tablvar}
\definecolor{darkgreen}{RGB}{0,100,0}
\providecommand{\red}{\color{red}}
\providecommand{\blue}{\color{blue}}
\providecommand{\green}{\color{green}}
\providecommand{\black}{\color{black}}
\providecommand{\white}{\color{white}}
\providecommand{\gray}{\color{gray}}
\providecommand{\cyan}{\color{cyan}}
\providecommand{\magenta}{\color{magenta}}
\providecommand{\yellow}{\color{yellow}}
\providecommand{\violet}{\color{violet}}
\providecommand{\orange}{\color{orange}}
\providecommand{\euro}{\texteuro}
\providecommand{\EUR}{\texteuro}
\providecommand{\ieme}{\textsuperscript{e}}
\providecommand{\degres}{\ensuremath{^\circ}}
\providecommand{\degree}{\ensuremath{^\circ}}
\providecommand{\vv}[1]{\overrightarrow{#1}}
\providecommand{\Coord}[1]{#1}
\providecommand{\decosix}{}
\providecommand{\starredbullet}{\ensuremath{\bullet}}
\pagestyle{empty}
\setlength{\parindent}{0pt}
\usepackage{adjustbox}\DeclareUnicodeCharacter{2236}{:}\DeclareUnicodeCharacter{2190}{\ensuremath{\leftarrow}}
\providecommand{\R}{\mathbb{R}}
\providecommand{\N}{\mathbb{N}}
\providecommand{\Z}{\mathbb{Z}}
\providecommand{\Q}{\mathbb{Q}}
\providecommand{\D}{\mathbb{D}}
\providecommand{\C}{\mathbb{C}}
\providecommand{\e}{{\mathrm{e}}}
\providecommand{\vect}[1]{\overrightarrow{#1}}
\providecommand{\vectt}[1]{\overrightarrow{\mathrm{#1}}}
\providecommand{\vvt}[1]{\overrightarrow{\mathrm{#1}}}
\providecommand{\np}[2][]{{#2~#1}}
\providecommand{\vv}[1]{{\overrightarrow{#1}}}
\renewcommand{\d}{{\mathrm{d}}}
\providecommand{\degres}{{^\circ}}
\providecommand{\degre}{{^\circ}}
\providecommand{\blue}{\color{blue}}
\providecommand{\red}{\color{red}}
\providecommand{\enskip}{\quad}
\providecommand{\llbracket}{\lbrack\!\lbrack}
\providecommand{\rrbracket}{\rbrack\!\rbrack}
\providecommand{\euro}{{\text{€}}}
\providecommand{\no}{{\text{n°}}}
\providecommand{\interval}[2]{{[#1;#2]}}
\providecommand{\intervFO}[2]{{[#1;#2[}}
\providecommand{\dotfill}{\ldots}

\providecommand{\ds}{\displaystyle}
\providecommand{\textding}[1]{#1}
\providecommand{\Oij}{\(\mathrm{O};\vec{\imath},\vec{\jmath}\)}
\providecommand{\Oijk}{\(\mathrm{O};\vec{\imath},\vec{\jmath},\vec{k}\)}
\providecommand{\Ouv}{\(\mathrm{O};\vec{u},\vec{v}\)}
\providecommand{\intervFF}[2]{\ensuremath{[#1;#2]}}
\renewcommand{\euro}{\ifmmode\text{\texteuro}\else\texteuro\fi}
\DeclareUnicodeCharacter{00B0}{\ifmmode^{\circ}\else\textdegree\fi}
\DeclareUnicodeCharacter{20AC}{\ifmmode\text{\texteuro}\else\texteuro\fi}
% Énoncés : Éducation nationale. Transcription et corrigés : APMEP, Denis Vergès et contributeurs. Découpage et classement : CoopMaths, Sébastien Lozano, Éric Elter et contributeurs. Adaptation de la présentation : Maths Explorer.
% CC BY-SA : https://coopmaths.fr/www/about/
% https://coopmaths.fr/alea/static/bac/2026/tex/bac_2026_05_sujet2_ameriquenord_1.tex
\begin{document}
\noindent\textit{APMEP / CoopMaths. CC BY-SA. Présentation adaptée par Maths Explorer.}\par\medskip
Un supermarché dispose d'un stock de tomates provenant de deux fournisseurs A et B.

II a été constaté que :

\begin{itemize}[label=$\bullet~$]
\item 91\,\% du stock de tomates est commercialisable ;
\item 60\,\% du stock de tomates provient du fournisseur A ;
\item parmi les tomates provenant du fournisseur A , la proportion de tomates commercialisables
est de 95\,\%.
\end{itemize}

On choisit au hasard une tomate dans le stock.

On désigne par :
\begin{itemize}[label=$\bullet~$]
\item $A$ l'évènement \og La tomate provient du fournisseur A \fg ;
\item $B$ l'évènement \og La tomate provient du fournisseur B \fg ;
\item $C$ l'évènement \og La tomate est commercialisable \fg.
\end{itemize}

Pour un évènement quelconque $E$, on note $P(E)$ la probabilité de $E$.

\medskip

\textbf{Partie A}

\medskip

\begin{enumerate}
\item Recopier l'arbre ci-dessous en complétant les pointillés.


\begin{center}
\pstree[treemode=R,levelsep=3.5cm, treesep=0.7cm]{\TR{}}
{\pstree{\TR{$A~$}\ncput*{\ldots}}
	{\TR{$C$} \ncput*{\ldots}
	\TR{$\overline{C}$} \ncput*{\ldots}
	}
\pstree{\TR{$B~$}\ncput*{\ldots}}
	{\TR{$C$} 
	\TR{$\overline{C}$} 
	}
}
\end{center}






















\item
	\begin{enumerate}
		\item Déterminer la probabilité que la tomate choisie soit commercialisable et provienne du fournisseur A
		\item Démontrer que $P_B(C) = 0,85$.
		\item La tomate choisie est non commercialisable. Le responsable des achats estime qu'il
y a deux fois moins de chance qu'elle provienne du fournisseur A que du fournisseur B. A-t-il raison ?
	\end{enumerate}
\end{enumerate}

\bigskip

\textbf{Partie B}

\medskip

On rappelle que 9\,\% des tomates du stock ne sont pas commercialisables.

\medskip

\begin{enumerate}
\item On prend $15$ tomates dans le stock au hasard et de manière indépendante. On considère que le stock est suffisamment important pour qu'on puisse assimiler ce prélèvement à un tirage aléatoire avec remise.

On note $X$ la variable aléatoire égale au nombre de tomates non commercialisables dans cet échantillon de $15$~tomates.
	\begin{enumerate}
		\item On admet que la variable aléatoire $X$ suit une loi binomiale. En préciser les paramètres.
		\item Déterminer la probabilité qu'exactement deux tomates soient non commercialisables.

\emph{On donnera la valeur arrondie au millième.}
		\item Déterminer la probabilité qu'au plus deux tomates soient non commercialisables.

\emph{On donnera la valeur arrondie au millième.}
	\end{enumerate}
\item On constitue désormais un échantillon de $n$ tomates, toujours dans les mêmes conditions, où $n$ désigne un entier naturel non nul.

On note $X_n$, la variable aléatoire égale au nombre de tomates non commercialisables et $F_n$ la variable aléatoire égale à la fréquence de tomates non commercialisables dans cet échantillon de $n$ tomates.

On a donc $F_n = \dfrac{X_n}{n}$.

On admet que la variable aléatoire $X_n$, suit la loi binomiale de paramètres $n$ et $0,09$.
	\begin{enumerate}
		\item Calculer l'espérance $E(F_n)$ et exprimer la variance $V(F_n)$ en fonction de $n$.
		\item Démontrer que $P(0,04 < F_n < 0,14) \geqslant 1 - \dfrac{32,76}{n}$.
		\item Le responsable des achats prélève dans le stock un échantillon de $500$ tomates. Il s'aperçoit que $55$ tomates ne sont pas commercialisables. 

Est-ce conforme à ce qu'il pouvait attendre ? Justifier la réponse.
	\end{enumerate}
\end{enumerate}

\medskip

\end{document}
